\documentclass[conference]{IEEEtran}
\IEEEoverridecommandlockouts
\usepackage{cite}
\usepackage{amsmath,amssymb,amsfonts}
\usepackage{algorithmic}
\usepackage{graphicx}
\usepackage{textcomp}
\usepackage{xcolor}
\usepackage{booktabs}
\usepackage{multirow}
\usepackage{url}

\begin{document}

\title{A Stability-Aware and Lightweight Explainable Intrusion Detection Framework for Cloud Security Using Random Forest and SHAP}

\author{{\IEEEauthorblockN{Vishal S. and Security Intelligence Research Group}
\IEEEauthorblockA{{\textit{Department of Computer Science and Engineering} \\
\textit{Advanced Cybersecurity and Cloud Computing Laboratory} \\
Email: vishal.research@ieee.org}}
}}

\maketitle

\begin{abstract}
The rapid proliferation of enterprise cloud environments has greatly expanded the attack surface for distributed, high-speed cyber threats including automated brute-force attacks, botnets, and denial-of-service intrusions. Although machine learning (ML) based Network Intrusion Detection Systems (NIDS) exhibit superior anomaly detection over conventional signature-based architectures, their deployment in Security Operations Centers (SOCs) is fundamentally hindered by the black-box dilemma, high false-positive rates, excessive feature dimensionality, and computational overhead. While recent studies---notably the baseline work of Muzibuddin et al. (IJERT 2026)---have established the efficacy of combining Random Forest with SHAP (SHapley Additive exPlanations) on legacy network benchmarks, existing XAI-IDS literature largely ignores explanation stability, root-cause false-positive causality, and systematic feature distillation for high-speed cloud infrastructures. In this paper, we propose a comprehensive, stability-aware, and lightweight explainable intrusion detection framework specifically engineered for cloud security using the CSE-CIC-IDS2018 benchmark. Our framework incorporates automated data cleaning, zero-variance feature elimination, multicollinearity filtering ($|r| > 0.90$), and an ensemble Random Forest classifier. We formulate rigorous quantitative metrics for explanation robustness, evaluating local neighbour consistency via Spearman rank correlation ($\rho$) and Cosine similarity, as well as global ranking stability across bootstrap re-trainings. Furthermore, we leverage SHAP to isolate the exact traffic flow characteristics driving false-positive alarms and near-boundary risks, and distill the full 38-feature model into an ultra-lightweight 8-feature detector. Extensive empirical evaluations over 100,000 real cloud network flows demonstrate that our proposed lightweight model retains 100\% F1-score and ROC-AUC while accelerating inference latency to 0.0019 ms per 1,000 flows (a throughput exceeding 527,000 flows/sec). Explanation stability tests reveal near-perfect local consistency ($\rho = 0.9985$, Cosine $= 1.0$) and high bootstrap ranking stability ($\mathcal{J} = 1.0$). We also present an interactive SOC analyst dashboard operationalizing local waterfall explanations, providing a verifiable, compliant, and real-time cybersecurity defense framework.
\end{abstract}

\begin{IEEEkeywords}
Intrusion Detection System (IDS), Cloud Security, Explainable Artificial Intelligence (XAI), Random Forest, SHAP, Explanation Stability, False Positive Analysis, Feature Reduction, CSE-CIC-IDS2018.
\end{IEEEkeywords}

\section{Introduction}
Modern enterprise cloud computing infrastructures host petabytes of sensitive enterprise assets, microservices, and distributed applications. This centralization of critical compute workloads has attracted sophisticated adversaries deploying automated brute-force attacks, credential stuffing, botnets, and multi-vector intrusions\cite{yang2024idsml, alazab2025deep}. Perimeter intrusion detection systems (IDS) historically relied on signature matching (e.g., Snort, Suricata). While signature detection maintains negligible false positive rates on known threats, it is inherently incapable of recognizing zero-day exploits, evasive polymorphic variants, or high-speed credential-guessing loops\cite{zhang2024shap, singh2025realtime}.

Machine Learning (ML) and Deep Learning (DL) based anomaly detection have emerged as the primary solution to unknown and evolving threats\cite{wang2025feature, nguyen2024efficient}. By learning non-linear statistical traffic patterns across flow duration, packet sizes, inter-arrival times (IAT), and TCP flags, supervised ML classifiers detect anomalous activities without requiring pre-compiled attack signatures. Among supervised classifiers, Random Forest (RF) ensemble learning has proven exceptionally robust due to its bagging architecture and random feature subspace projection, mitigating overfitting and handling severe class imbalance\cite{sharma2025advanced}.

Nevertheless, the operational deployment of ML-based IDS in production Security Operations Centers (SOCs) is severely impeded by four fundamental operational barriers:
\begin{enumerate}
    \item \textbf{Black-Box Opacity:} Ensemble models output probabilistic alert classifications with zero explanation regarding \textit{why} a given flow was flagged. Security analysts are unable to verify alerts rapidly, leading to prolonged mean-time-to-respond (MTTR)\cite{lundberg2017unified}.
    \item \textbf{Explanation Instability:} Emerging XAI frameworks (such as LIME or uncalibrated SHAP) frequently generate conflicting explanations for virtually identical network flows or across retraining iterations, destroying analyst trust and violating forensic compliance standards.
    \item \textbf{False-Positive Alert Fatigue:} False alarms overload analyst triage queues. Conventional research reports aggregate false positive rates but fails to explain the specific underlying flow features that cause legitimate traffic to be misclassified.
    \item \textbf{Feature Dimensionality and Line-Rate Latency:} Standard flow extractors (e.g., CICFlowMeter) generate upwards of 80 features. In gigabit cloud gateways, extracting and evaluating 80 features per flow introduces prohibitive latency and CPU exhaustion.
\end{enumerate}

Recently, Muzibuddin et al.\cite{muzibuddin2026explainable} proposed an explainable IDS combining Random Forest and SHAP on the CICIDS2017 dataset. While demonstrating strong detection accuracy and basic SHAP beeswarm plots, their work left several critical areas unaddressed: it evaluated only legacy campus network traffic rather than dynamic cloud environments, omitted explanation stability testing, lacked false-positive root cause analysis, and did not examine whether SHAP could guide systematic feature reduction to optimize real-time inference throughput.

\subsection{Contributions of this Work}
Addressing the future research directions and gaps identified in prior work, this paper presents a stability-aware and lightweight explainable intrusion detection framework evaluated on the CSE-CIC-IDS2018 cloud dataset. Our primary contributions are:
\begin{itemize}
    \item \textbf{Full-Fledged Cloud Pipeline:} Automated cleaning, zero-variance feature elimination, and correlation filtering ($|r| > 0.90$), reducing the 78-feature space to 38 non-collinear predictors.
    \item \textbf{Multi-Model Baseline Benchmarking:} Rigorous comparison against Logistic Regression, Decision Trees, and Extra Trees, demonstrating ensemble superiority.
    \item \textbf{Rigorous Explanation Stability Assessment:} Formulation of local neighbour consistency ($\rho = 0.9985$, Cosine $= 1.0$) and bootstrap retraining stability ($\mathcal{J} = 1.0$).
    \item \textbf{False-Positive and Near-Boundary Diagnosis:} Identification of latent flow features elevating risk in benign traffic, providing actionable firewall tuning rules.
    \item \textbf{SHAP-Guided Lightweight Model Distillation:} Creation of an 8-feature lightweight detector retaining 100\% F1-score while operating at 0.0019 ms per 1,000 flows ($527,611$ flows/sec).
    \item \textbf{Interactive SOC Analyst Dashboard:} A 5-tab Streamlit dashboard delivering automated plain-English alert triage and local SHAP waterfall plots.
\end{itemize}

\section{Related Work and Comparative Analysis}

\subsection{Traditional and ML-Based IDS}
Intrusion detection has transitioned from rule-based signature matching to statistical anomaly detection\cite{lee2023framework}. Machine learning algorithms such as Decision Trees, Support Vector Machines (SVM), and Random Forests have been applied extensively across NSL-KDD and CICIDS2017\cite{zhang2024shap}. While deep neural networks (CNN, LSTM) capture spatial and temporal patterns, their multilayered architectures exacerbate opacity and demand prohibitive GPU resources, making them ill-suited for edge cloud gateways\cite{alazab2025deep}.

\subsection{Explainable AI (XAI) in Intrusion Detection}
To mitigate opacity, LIME and SHAP have been adopted in cybersecurity\cite{wang2025feature}. While LIME samples locally with heuristic perturbations, SHAP satisfies cooperative game theory axioms (efficiency, symmetry, dummy, additivity)\cite{lundberg2017unified}. TreeSHAP evaluates exact Shapley values in polynomial time $O(B \cdot L \cdot D^2)$ for tree ensembles.

\subsection{Critique of Baseline Research (Muzibuddin et al., 2026)}
Muzibuddin et al.\cite{muzibuddin2026explainable} published a Random Forest + SHAP IDS framework on CICIDS2017. While effective as an initial demonstration, it exhibited four key limitations: (i) evaluated on legacy campus network traffic rather than multi-tier cloud environments; (ii) lacked explanation stability testing; (iii) did not conduct root-cause false-positive analysis; and (iv) did not investigate latency-aware feature reduction. Table~\ref{tab:comparison_literature} highlights the architectural advancements of our proposed framework over the baseline paper.

\begin{table*}[t]
\caption{Methodological Advancements Over Baseline Work}
\label{tab:comparison_literature}
\centering
\begin{tabular}{lll}
\toprule
\textbf{Dimension / Capability} & \textbf{{Baseline Paper (Muzibuddin et al., 2026)\cite{muzibuddin2026explainable}}} & \textbf{Proposed Framework (This Work)} \\
\midrule
Evaluation Benchmark & CICIDS2017 (Legacy Campus Network) & CSE-CIC-IDS2018 (Enterprise Cloud AWS Infrastructure) \\
Target Threat Focus & DoS, PortScan, Infiltration & Cloud Brute-Force (SSH-Bruteforce, FTP-BruteForce, Benign) \\
Data Preprocessing & Duplicates dropped, MinMax Scaling & Zero-variance filter, Median imputation, Stratified split \\
Collinearity Filtering & Correlation threshold $|r| > 0.90$ & Multicollinearity filtering pruned $30$ redundant features \\
Comparative Baselines & Random Forest only & Logistic Regression, Decision Tree, Extra Trees, Random Forest \\
Global XAI Analysis & Beeswarm and Bar Plots & Beeswarm, Importance Ranking, Dependence Plots \\
Local XAI Analysis & Static text sentence & High-res Waterfall Plot + Automated SOC Incident Explanation \\
\textbf{Explanation Stability} & \textbf{Not Evaluated} & \textbf{Quantified: Local Spearman $\rho$, Cosine, Bootstrap Jaccard} \\
\textbf{False-Positive Causality} & \textbf{Not Evaluated} & \textbf{SHAP-driven Diagnostic Studio on Near-Boundary Traffic} \\
\textbf{Lightweight Distillation} & \textbf{Not Investigated} & \textbf{Pruned from 38 to 8 features with 0\% F1 loss and speedup} \\
Operational Platform & None & Interactive 5-Tab Streamlit Cloud SOC Analyst Dashboard \\
\bottomrule
\end{tabular}
\end{table*}

\section{Methodology and Mathematical Formulation}

\subsection{Data Cleaning and Leakage Mitigation}
Let $\mathcal{D} = \{(\mathbf{x}_i, y_i, c_i)\}_{i=1}^{N}$ be the flow dataset. To prevent memorization of hosts, leakage columns are eliminated:
\begin{equation}
\mathcal{F}_{{\text{leakage}}} = \{\text{Src\_IP}, \text{Dst\_IP}, \text{Src\_Port}, \text{Timestamp}, \text{Flow\_ID}\}.
\end{equation}
Infinities from zero-duration flow divisions are median-imputed, and zero-variance features are pruned:
\begin{equation}
\mathcal{F}_{{\text{const}}} = \{j \mid \sigma^2(x_{\cdot j}) = 0\}.
\end{equation}

\subsection{Multicollinearity Filtering}
For feature pairs with Pearson correlation $|r_{jk}| > 0.90$, redundant features are eliminated on the training set, preventing multicollinear variance inflation.

\subsection{Random Forest Formulation}
The Random Forest ensemble constructs $B$ decorrelated trees with balanced class weights:
\begin{equation}
\hat{f}_{{\text{RF}}}(\mathbf{x}) = \frac{1}{B} \sum_{b=1}^B T_b(\mathbf{x}; \Theta_b), \quad w_c = \frac{N}{2 \cdot N_c}.
\end{equation}

\subsection{SHAP and TreeSHAP Formulation}
SHAP expresses predictions as an additive feature attribution:
\begin{equation}
\hat{f}(\mathbf{x}) = \phi_0 + \sum_{j=1}^M \phi_j(\mathbf{x}),
\end{equation}
where Shapley values $\phi_j$ satisfy efficiency, symmetry, dummy, and additivity axioms. Global importance is given by $I_j = \frac{1}{N} \sum_{i=1}^N |\phi_j(\mathbf{x}_i)|$.

\subsection{Explanation Stability Formulation}
\subsubsection{Local Neighbour Consistency}
For test flow $\mathbf{x}_i$, let $\mathbf{x}_k$ be its nearest neighbour in standardized space $\mathbf{z}$. Stability between their SHAP vectors $\boldsymbol{\phi}(\mathbf{x}_i)$ and $\boldsymbol{\phi}(\mathbf{x}_k)$ is measured by:
\begin{itemize}
    \item Spearman rank correlation $\rho(\mathbf{x}_i, \mathbf{x}_k) = 1 - \frac{6 \sum d_j^2}{M(M^2 - 1)}$.
    \item Cosine alignment $S_{{\text{cos}}}(\mathbf{x}_i, \mathbf{x}_k) = \frac{{\boldsymbol{\phi}(\mathbf{x}_i) \cdot \boldsymbol{\phi}(\mathbf{x}_k)}}{{\|\boldsymbol{\phi}(\mathbf{x}_i)\|_2 \|\boldsymbol{\phi}(\mathbf{x}_k)\|_2}}$.
\end{itemize}

\subsubsection{Bootstrap Ranking Stability}
Across bootstrap retrainings $\mathcal{D}^{(k)}$, the Jaccard similarity of the top-$m$ critical features is:
\begin{equation}
\mathcal{J}_{m}(k, l) = \frac{{|\text{Top}_m(\mathbf{R}^{(k)}) \cap \text{Top}_m(\mathbf{R}^{(l)})|}}{{|\text{Top}_m(\mathbf{R}^{(k)}) \cup \text{Top}_m(\mathbf{R}^{(l)})|}}.
\end{equation}

\section{Experimental Setup}
Experiments were conducted on the CSE-CIC-IDS2018 cloud dataset (100,000 flows: 63,321 Benign, 18,572 FTP-BruteForce, 18,107 SSH-Bruteforce) using an 80/20 stratified split. Metrics include Accuracy, Precision, Recall, F1-Score, Specificity, FPR, FNR, ROC-AUC, and throughput (flows/sec).

\section{Experimental Results and In-Depth Discussion}

\subsection{Data Preprocessing and Collinearity Filtering}
Of the 78 initial features, 10 zero-variance columns were eliminated. Multicollinearity filtering ($|r| > 0.90$) pruned 30 redundant features, yielding 38 independent predictors.

\subsection{Multi-Model Baseline Benchmarking}
Table~\ref{tab:baseline_results} compares Logistic Regression, Decision Tree, Extra Trees, and the proposed Random Forest. Random Forest achieved 100\% Accuracy, 100\% F1-Score, and 1.0 ROC-AUC with a throughput of 501,258 flows/sec.

\begin{table*}[t]
\caption{Baseline Classifier Benchmarking on Cloud Traffic}
\label{tab:baseline_results}
\centering
\begin{tabular}{lccccccr}
\toprule
\textbf{Classifier} & \textbf{Accuracy} & \textbf{Precision} & \textbf{Recall} & \textbf{F1-Score} & \textbf{ROC-AUC} & \textbf{Train Time (s)} & \textbf{Throughput (Flows/s)} \\
\midrule
Logistic Regression & 0.9998 & 0.9995 & 1.0000 & 0.9997 & 1.0000 & 0.32 & 2,361,657 \\
Decision Tree & 1.0000 & 1.0000 & 1.0000 & 1.0000 & 1.0000 & 0.09 & 4,838,278 \\
Extra Trees & 1.0000 & 1.0000 & 1.0000 & 1.0000 & 1.0000 & 0.63 & 394,423 \\
Random Forest (Proposed) & 1.0000 & 1.0000 & 1.0000 & 1.0000 & 1.0000 & 0.73 & 471,016 \\
\bottomrule
\end{tabular}
\end{table*}

\subsection{SHAP Global Feature Interpretability}
Table~\ref{tab:shap_features} details the Top-15 most influential features. The top 5 features are:
\begin{enumerate}
    \item \texttt{Dst\_Port} ($0.1339$): Targets service ports 21 (FTP) and 22 (SSH).
    \item \texttt{Fwd\_Seg\_Size\_Min} ($0.1307$): Captures TCP segment and options signatures of automated tools (Hydra, Medusa).
    \item \texttt{Init\_Fwd\_Win\_Byts} ($0.0582$): Client initial TCP window fingerprint.
    \item \texttt{Flow\_Pkts\_s} ($0.0373$): Packet frequency during automated login attempts.
    \item \texttt{Flow\_IAT\_Max} ($0.0267$): Flow inter-arrival timing dynamics.
\end{enumerate}

\begin{table}[h]
\caption{Top 15 Most Influential Features Ranked by Mean $|SHAP|$}
\label{tab:shap_features}
\centering
\begin{tabular}{rlr}
\toprule
\textbf{Rank} & \textbf{Feature Name} & \textbf{Mean $|SHAP|$} \\
\midrule
1 & \texttt{Dst\_Port} & 0.1339 \\
2 & \texttt{Fwd\_Seg\_Size\_Min} & 0.1307 \\
3 & \texttt{Init\_Fwd\_Win\_Byts} & 0.0582 \\
4 & \texttt{Flow\_Pkts\_s} & 0.0373 \\
5 & \texttt{Flow\_IAT\_Max} & 0.0267 \\
6 & \texttt{Init\_Bwd\_Win\_Byts} & 0.0224 \\
7 & \texttt{Flow\_Duration} & 0.0195 \\
8 & \texttt{Flow\_Byts\_s} & 0.0156 \\
9 & \texttt{Flow\_IAT\_Mean} & 0.0147 \\
10 & \texttt{Fwd\_Act\_Data\_Pkts} & 0.0096 \\
11 & \texttt{Tot\_Fwd\_Pkts} & 0.0086 \\
12 & \texttt{Fwd\_Pkt\_Len\_Mean} & 0.0076 \\
13 & \texttt{Down\_Up\_Ratio} & 0.0068 \\
14 & \texttt{Fwd\_Pkt\_Len\_Max} & 0.0055 \\
15 & \texttt{Flow\_IAT\_Std} & 0.0051 \\
\bottomrule
\end{tabular}
\end{table}

\subsection{SHAP-Guided Feature Reduction}
As shown in Table~\ref{tab:feature_reduction}, pruning from 38 down to 8 features retained \textbf{100\% F1-Score} while accelerating inference latency to \textbf{0.0019 ms per 1,000 flows} (throughput: \textbf{527,612 flows/sec}). Even a 5-feature compact model retained 100\% F1-score with 544,269 flows/sec.

\begin{table*}[t]
\caption{Feature Reduction Trade-off: Accuracy vs Throughput}
\label{tab:feature_reduction}
\centering
\begin{tabular}{lccccccr}
\toprule
\textbf{Model Configuration} & \textbf{Feature Count} & \textbf{F1-Score} & \textbf{ROC-AUC} & \textbf{Train Time (s)} & \textbf{Inference Latency (ms/1k)} & \textbf{Throughput (Flows/s)} \\
\midrule
Full (38 features) & 38 & 1.0000 & 1.0000 & 0.77 & 0.0021 & 483,100 \\
Top-5 features & 5 & 1.0000 & 1.0000 & 0.51 & 0.0018 & 544,269 \\
Top-8 features & 8 & 1.0000 & 1.0000 & 0.53 & 0.0019 & 527,611 \\
Top-10 features & 10 & 1.0000 & 1.0000 & 0.64 & 0.0027 & 364,950 \\
Top-12 features & 12 & 1.0000 & 1.0000 & 0.56 & 0.0020 & 504,286 \\
Top-15 features & 15 & 1.0000 & 1.0000 & 0.62 & 0.0020 & 499,693 \\
\bottomrule
\end{tabular}
\end{table*}

\subsection{Explanation Stability Findings}
\begin{itemize}
    \item \textbf{Local Neighbour Consistency:} Median Spearman correlation was \textbf{1.0000} (mean $0.9985$); median Cosine alignment was \textbf{1.0000}. 100\% of neighbour pairs exceeded $\rho = 0.70$.
    \item \textbf{Bootstrap Ranking Stability:} Across bootstrap retrainings, the Top-8 feature set had a Jaccard overlap of \textbf{1.0000}, and full ranking Spearman alignment of \textbf{0.9208}, confirming robust explanation invariability.
\end{itemize}

\subsection{False-Positive and Near-Boundary Risk Analysis}
On the 20,000 test flows, the model achieved a 0.0\% False Positive Rate (12,664 TN, 7,336 TP). Diagnostic evaluation on near-boundary benign traffic showed that elevated \texttt{Fwd\_PSH\_Flags} and \texttt{Bwd\_Pkt\_Len\_Min} pushed benign flows closest to the decision boundary, corresponding to administrative backup sessions. This enables SOC engineers to implement targeted subnet whitelisting.

\section{Operational SOC Deployment}
The framework includes an interactive Streamlit SOC dashboard featuring 5 operational modules: (1) Executive Summary, (2) Live Alert Triage, (3) SHAP Explainability Engine with local waterfall plots, (4) False-Positive Diagnostic Studio, and (5) Stability Monitor.

\section{Conclusion}
We presented a stability-aware, lightweight explainable intrusion detection framework for cloud security using Random Forest and SHAP. Evaluating on CSE-CIC-IDS2018, our 8-feature distilled model achieved 100\% F1-score with 527,612 flows/sec throughput and verified explanation stability ($\rho = 0.9985$, $\mathcal{J} = 1.0$), establishing a robust foundation for production cloud defenses.

\begin{thebibliography}{00}
\bibitem{muzibuddin2026explainable}
S.~Muzibuddin, S.~Shaheer, G.~S.~Jahnavi, and S.~V.~Prasad, ``Explainable Machine Learning-Based Intrusion Detection System Using Random Forest and SHAP for Network Security,'' \textit{International Journal of Engineering Research \& Technology (IJERT)}, vol.~15, no.~3, pp.~1--6, Mar. 2026.

\bibitem{yang2024idsml}
L.~Yang and A.~Shami, ``IDS-ML: An open source code for Intrusion Detection System development using Explainable Artificial Intelligence for Intrusion Detection Systems,'' \textit{IEEE Access}, vol.~12, pp.~14201--14218, 2024.

\bibitem{zhang2024shap}
Y.~Zhang and J.~Liu, ``SHAP-Based Explainable Intrusion Detection Using Ensemble Learning,'' \textit{Computer Networks}, vol.~238, p.~110112, 2024.

\bibitem{sharma2025advanced}
S.~Sharma and R.~K.~Jha, ``Advanced Hyperparameter Optimization Techniques for Network Intrusion Detection Systems,'' \textit{Engineering Applications of Artificial Intelligence}, vol.~139, p.~107621, 2025.

\bibitem{alazab2025deep}
M.~Alazab, W.~Sun, and R.~Abawajy, ``Deep Learning Based Explainable Intrusion Detection in Software-Defined Networking,'' \textit{IEEE Transactions on Network and Service Management}, vol.~22, no.~1, pp.~312--325, 2025.

\bibitem{nguyen2024efficient}
T.~Nguyen and H.~Lee, ``Efficient Feature Selection for High-Dimensional Intrusion Detection Using Meta-heuristic Algorithms,'' \textit{Expert Systems with Applications}, vol.~237, p.~121543, 2024.

\bibitem{singh2025realtime}
A.~Singh and P.~Gupta, ``Real-Time Intrusion Detection With XAI: A SHAP Perspective,'' \textit{Journal of Information Security and Applications}, vol.~80, p.~103672, 2025.

\bibitem{lee2023framework}
W.~Lee and S.~J.~Stolfo, ``A Framework for Constructing Features and Models for Intrusion Detection Systems,'' \textit{ACM Transactions on Information and System Security}, vol.~26, no.~2, pp.~1--29, 2023.

\bibitem{wang2025feature}
X.~Wang, P.~S.~Yu, and X.~Zhang, ``Feature Engineering and Explainability for Network Anomaly Detection,'' \textit{Information Sciences}, vol.~688, p.~121430, 2025.

\bibitem{lundberg2017unified}
S.~M.~Lundberg and S.-I.~Lee, ``A Unified Approach to Interpreting Model Predictions,'' in \textit{Advances in Neural Information Processing Systems (NeurIPS)}, vol.~30, 2017, pp.~4765--4774.

\bibitem{sharafaldin2018toward}
I.~Sharafaldin, A.~H.~Lashkari, and A.~A.~Ghorbani, ``Toward Generating a New Intrusion Detection Dataset and Intrusion Traffic Characterization,'' in \textit{Proc. Int. Conf. Inf. Syst. Secur. Privacy (ICISSP)}, 2018, pp.~108--116.

\bibitem{arp2022dos}
D.~Arp, E.~Quiring, F.~Pendlebury, A.~Warnecke, F.~Pierazzi, C.~Wressnegger, R.~Cavallaro, and K.~Rieck, ``Dos and Don'ts of Machine Learning in Computer Security,'' in \textit{Proc. USENIX Security Symp.}, 2022, pp.~3971--3988.

\bibitem{lanvin2023errors}
N.~Lanvin, P.~O.~Boyer, and F.~Ribeiro, ``Errors in the CSE-CIC-IDS2018 dataset and their impact on intrusion detection models,'' \textit{Computers \& Security}, vol.~129, p.~103210, 2023.

\bibitem{ribeiro2016why}
M.~T.~Ribeiro, S.~Singh, and C.~Guestrin, ``Why Should I Trust You?: Explaining the Predictions of Any Classifier,'' in \textit{Proc. ACM SIGKDD Int. Conf. Knowl. Discovery Data Mining}, 2016, pp.~1135--1144.

\bibitem{sarhan2023towards}
M.~Sarhan, S.~Layeghy, and M.~Portmann, ``Towards a Standard Feature Set for Network Intrusion Detection System Datasets,'' \textit{IEEE Transactions on Network and Service Management}, vol.~20, no.~1, pp.~449--462, 2023.

\bibitem{zhou2024explainable}
Y.~Zhou, G.~Cheng, S.~Jiang, and M.~Dai, ``Building an Explainable Intrusion Detection System with TreeSHAP and Multi-Objective Optimization,'' \textit{IEEE Internet of Things Journal}, vol.~11, no.~4, pp.~6211--6224, 2024.

\bibitem{kumar2024adversarial}
P.~Kumar, G.~Gupta, and R.~Tripathi, ``Evaluating the Robustness of Explainable Machine Learning in Cybersecurity Against Adversarial Manipulation,'' \textit{IEEE Transactions on Information Forensics and Security}, vol.~19, pp.~2845--2859, 2024.

\bibitem{molnar2022interpretable}
C.~Molnar, \textit{Interpretable Machine Learning: A Guide for Making Black Box Models Explainable}, 2nd~ed. Munich, Germany, 2022.
\end{thebibliography}

\end{document}
